% Partie : méthodes
% Chapitre : induction mathématique
% Section : récurrence sur les nombres naturels

\documentclass[../../../include/open-logic-section]{subfiles}

\begin{document}
\olfileid[fr]{mth}{ind}{inN}

\olsection{Récurrence sur~\texorpdfstring{$\Nat$}{ℕ}}

La récurrence sur les entiers naturels permet d'établir
qu'une propriété vaut pour chacun d'eux. Elle utilise
le fait qu'en partant de~$0$
et en ajoutant~$1$ de façon
répétée, on atteint finalement
chaque nombre naturel.
Pour montrer qu'une propriété
vaut pour tous ces nombres,
on peut donc
(1)~établir qu'elle vaut
pour~$0$, puis
(2)~montrer que, chaque fois
qu'elle vaut pour un nombre~$n$,
elle vaut aussi pour le
suivant, $n+1$.
Abrégeons «~le nombre~$n$
possède la propriété~$P$~»
par $P(n)$, et de même
pour $P(k)$, etc.
Une démonstration par
récurrence de $P(n)$
pour tout $n \in \Nat$
comprend alors :
\begin{enumerate}
\item une démonstration de $P(0)$ ;
\item une démonstration que,
  pour tout~$k$, si $P(k)$,
  alors $P(k+1)$.
\end{enumerate}
Examinons précisément pourquoi
ces deux étapes suffisent.
La première donne~$P(0)$.
La seconde, appliquée avec
$k=0$, donne
«~si $P(0)$, alors $P(0+1)$~»,
c'est-à-dire
«~si $P(0)$, alors $P(1)$~».
Par modus ponens, nous
obtenons~$P(1)$.
En appliquant de nouveau
la seconde étape, cette
fois avec $k=1$,
nous obtenons
«~si $P(1)$, alors $P(2)$~».
Puisque nous venons
d'établir $P(1)$,
un autre modus ponens
donne~$P(2)$.
Et ainsi de suite :
pour tout~$n$, après
avoir répété l'étape
$n$~fois, nous obtenons
finalement~$P(n)$.
Les deux étapes établissent
donc $P(n)$ pour tout
$n \in \Nat$.

Prenons un exemple.
Combien de sommes différentes
peut-on obtenir en lançant
$n$ dés ? Commençons,
même si cela paraît étrange,
avec $0$ dé. Sans dés,
il n'y a qu'une somme
possible : aucun point,
soit~$0$. Le nombre
de sommes possibles est
donc~$1$. Avec un seul
dé, $n=1$, on obtient
six valeurs, de $1$ à~$6$.
Avec deux dés, toutes
les sommes de $2$ à
$12$ sont possibles :
il y en a~$11$.
Avec trois dés, on
peut obtenir toutes
les sommes de $3$
à $18$, soit~$16$.
La suite $1$, $6$,
$11$, $16$ suggère
une règle : la réponse
serait-elle $5n+1$ ?
C'est en tout cas le
nombre maximal de
valeurs distinctes
possibles, car il y
a précisément $5n+1$
entiers entre $n$,
la plus petite somme
possible avec $n$
dés (tous montrant~$1$),
et $6n$, la plus
grande (tous
montrant~$6$).

\begin{thm}
  Avec $n$ dés, on peut
  obtenir chacune des
  $5n+1$ valeurs
  comprises entre
  $n$ et $6n$.
\end{thm}

\begin{proof}
Notons $P(n)$ l'énoncé :
«~Avec $n$ dés,
on peut obtenir
tout entier compris
entre $n$ et~$6n$.~»
Le cas sans dé a été établi ci-dessus.
Pour démontrer cet énoncé
par récurrence, établissons :
%
%
%
%
%
%
%
%
%
%
\begin{enumerate}
\item la \emph{base de la récurrence}
  $P(1)$ : avec un seul dé,
  toutes les sommes de
  $1$ à $6$ sont possibles ;
\item l'\emph{étape de récurrence} :
  pour tout $k$,
  si $P(k)$, alors~$P(k+1)$.
\end{enumerate}

Pour (1), il suffit d'examiner
un dé à six faces : chacune
des valeurs de $1$ à~$6$
figure sur l'une de ses faces.
Un seul dé permet donc
d'obtenir toutes les valeurs
de cet intervalle.

Pour (2), supposons
l'antécédent de l'implication,
$P(k)$. Cette supposition
s'appelle l'\emph{hypothèse
de récurrence}.
Nous l'utiliserons pour
démontrer~$P(k+1)$.
La difficulté consiste
à exprimer les sommes
obtenues avec $k+1$ dés
à partir de celles que
l'on peut obtenir avec
$k$ dés et de la valeur
du dé supplémentaire :
c'est ainsi que nous
pourrons exploiter
l'hypothèse de récurrence.

Celle-ci dit qu'avec
$k$ dés, on peut obtenir
tous les entiers de
$k$ à~$6k$. Si le
$(k+1)$-ième dé montre~$1$,
on ajoute~$1$ à la
somme : toutes les
valeurs de $k+1$
à $6k+1$ sont donc
possibles. Il reste
les valeurs de $6k+2$
à~$6k+6$. Pour les obtenir, il suffit d'avoir~$6$ sur
chacun des $k$ premiers dés et la valeur voulue,
comprise entre $2$ et~$6$, sur le dernier.
Avec $k+1$ dés,
on peut ainsi obtenir
tous les entiers de
$k+1$ à $6(k+1)$ :
nous avons démontré
$P(k+1)$ à partir
de l'hypothèse~$P(k)$.
\end{proof}

Très souvent, on recourt
à la récurrence pour
démontrer un résultat
concernant une suite
d'objets (nombres,
ensembles, etc.) qui
est elle-même définie
«~par récurrence~» :
le $(n+1)$-ième objet
est défini à partir
du $n$-ième.
Par exemple, la
somme~$s_n$ des
nombres naturels
jusqu'à~$n$ peut
être définie par
\begin{align*}
  s_0 & = 0\\
  s_{n+1} & = s_n + (n+1).
\end{align*}
Cette définition donne :
\begin{align*}
  s_0 & = 0,\\
  s_1 & = s_0 + 1 && = 1,\\
  s_2 & = s_1 + 2 && = 1 + 2 = 3,\\
  s_3 & = s_2 + 3 && = 1 + 2 + 3 = 6,
  \text{ etc.}
\end{align*}
Nous pouvons maintenant
démontrer par récurrence
que $s_n=n(n+1)/2$.

\begin{prop}
  $s_n = n(n+1)/2$.
\end{prop}

\begin{proof}
  Il faut démontrer
  (1) $s_0=0\cdot(0+1)/2$
  et
  (2) si
  $s_k=k(k+1)/2$,
  alors
  $s_{k+1}=(k+1)(k+2)/2$.
  L'assertion (1)
  est immédiate.
  Pour (2), supposons
  l'hypothèse de récurrence
  $s_k=k(k+1)/2$.
  Nous devons en déduire
  $s_{k+1}=(k+1)(k+2)/2$.

  Que vaut $s_{k+1}$ ?
  Par définition,
  $s_{k+1}=s_k+(k+1)$.
  En substituant
  $s_k=k(k+1)/2$
  et en calculant,
  nous obtenons :
  \begin{align*}
    s_{k+1}
      & = \frac{k(k+1)}{2}+(k+1)\\
      & = \frac{k(k+1)}{2}
        +\frac{2(k+1)}{2}\\
      & = \frac{k(k+1)+2(k+1)}{2}\\
      & = \frac{(k+2)(k+1)}{2}.
  \end{align*}
\end{proof}

La leçon importante est
la suivante : lorsqu'on
étudie une suite~$a_n$
définie par récurrence,
une démonstration par
récurrence est la
méthode naturelle.
Même si la suite des
objets n'est pas
définie de cette
manière, comme dans
l'exemple des dés,
on peut utiliser
la récurrence dès
qu'on sait relier
le cas~$k+1$ au
cas~$k$.

\end{document}
